\documentclass[12pt]{article}

% Packages
\usepackage[a4paper, margin=1in]{geometry} % Adjust margins
\usepackage{setspace} % For line spacing
\usepackage{tocloft} % Optional: better spacing for TOC
\usepackage[backend=biber,style=apa]{biblatex} % Load biblatex with APA style
\addbibresource{references.bib} % Link to your .bib file

% Set line spacing
\onehalfspacing % or use \doublespacing for double spacing

% Title information
\title{\textbf{\large{Preferences on Tax Enforcement and Tax Administration}}}
\author{\normalsize{Tim Kircali, Stefanie Stantcheva, Matthias Weber}}
\date{\normalsize{30.10.2025}}
\begin{document}
\maketitle
\pagenumbering{arabic}
%\begin{center}
%    \textbf{Abstract}
%\end{center}
%This study explores public preferences regarding tax enforcement and tax administration, focusing on attitudes toward information sharing and government-provided tax filing services. Using a planned survey on a representative sample of the U.S. population, we examine views on comprehensive third-party reporting, international exchange of information, and government-provided tax filing services. We explore underlying considerations such as data privacy, tax evasion, perceptions of fairness and inequality, and tax filing costs. To identify causal drivers of policy preferences, we include educational video treatments that provide targeted information on key considerations. This experimental approach helps to reveal how specific concerns shape support for tax enforcement and tax administration, offering insights relevant for future reforms.


%This study explores public preferences regarding tax enforcement and tax administration, focusing on attitudes towards information-sharing and government-provided tax filing services. Through a planned survey on a representative sample of the U.S.\ population, we examine views on comprehensive third-party reporting, international exchange of information and government provided tax filing services. We ask questions to explore important underlying considerations such as data privacy, tax evasion (extent + distribution) and tax filing costs. To identify causal drivers of policy preferences, we include video treatments that provide targeted information on important considerations. This experimental approach helps to reveal how specific concerns shape support for tax enforcement and tax administration, insights which are important for future reforms.




%This study explores public preferences regarding tax enforcement and tax administration, focusing on attitudes towards the expansion of third-party reporting and the role of governments in providing tax software and prepopulated tax returns. We conduct a survey experiment on a representative sample of the U.S.\ population to analyse perceptions, beliefs and important considerations underlying 

%public preferences regarding tax enforcement and tax administration, focusing on attitudes towards information-sharing and free We explore attitudes and considerations towards comprehensive third-party reporting and international
%automatic exchange of information, and toward free tax software and prepopulated tax returns. 
%tax filing service provision. comprehensive third-party reporting, filing services and underlying considerations such as privacy, tax evasion, views on fairness and filing costs.


%We vary the information that participants obtain to make causal inference on the role that information on tax evasion, inequality, and tax filing costs plays in views on comprehensive third-party reporting. Our results show that without additional information, about half of the respondents are in favor of comprehensive third-party reporting. This number increases up to almost three quarters for participants in our full information treatmet. Free tax software and the pre-population of tax returns by the tax authority are supported by a majority of respondents who have not received additional information and by a vast majority of respondents in our information treatments.


%\begin{center}
%\textbf{Abstract} 
%\end{center}
%This study explores public preferences regarding tax enforcement and tax administration, focusing on attitudes towards the expansion of third-party reporting and the role of governments in providing tax software and prepopulated tax returns. Through a planned survey in the US, Germany, Switzerland and Denmark, we examine views on comprehensive third-party reporting, filing services and underlying considerations such as privacy, tax evasion, views on fairness and filing costs. To identify causal drivers of policy preferences, we include video treatments that provide targeted information on each consideration. This experimental approach helps to reveal how specific concerns shape support for tax enforcement and administrative reforms.
\begin{center}
\textbf{Extended Abstract}  
\end{center}
For governments to fulfill their responsibilities, such as providing security, infrastructure, education, health, and social protection, it is essential that taxes that are due are actually paid. Nevertheless, tax evasion remains a serious problem in many countries. Research has shown that automatic exchange of information (AEoI) and third-party reporting are highly effective tools in reducing tax evasion. Moreover, comprehensive third-party reporting not only helps enforce tax laws, but also improves the services that tax administrations can provide to citizens. In some countries, taxpayers face high compliance costs and are urged to purchase commercial tax software to file tax returns. In others, governments ease the burden of tax filing by offering free tax software and pre-populated tax returns. Despite the clear benefits of AEoI and third-party reporting, many countries are still hesitant or reluctant to adopt these tools more widely.

This study explores citizens' preferences for how taxes should be enforced and what types of services tax administrations should provide. In particular, we examine attitudes towards expanding third-party reporting and offering government-provided tax filing services. In a planned survey among residents of the US, Germany, Switzerland and Denmark, we ask: Should governments collect more information through third parties? What types of data and which taxpayers should be included? And, if we switch to the tax administrations' service perspective, should tax authorities reduce filing costs by providing free software and prepopulated tax returns? These questions are closely related, as the provision of prepopulated tax returns requires automatic access to taxpayer data.

Furthermore, we try to understand how people reason about these views? What are the considerations driving support for comprehensive third party reporting and what is the role of (mis)information? There might be a big number of considerations underlying support for tax enforcement policies, however, we try to understand the role of four economically convincing factors underlying policy support; data privacy concerns, considerations of tax evasion and tax revenue, considerations regarding fairness \& inequality and tax filing concerns. In addition, previous research suggests that political leanings \& government views might play a strong role in the formation of the policy preferences. To shed more light on how these underlying considerations link to policy views, we include questions to elicit them. 

Finally, to understand the role of (mis)information and to establish a causal link between considerations and policy views, we experimentally provide respondents with informational video treatments. Four video treatments each specifically address one of the four underlying considerations mentioned above in isolation, while one video treatment, the full-info treatment, includes all information. This allows us to observe how the considerations of the treated groups change and, more interestingly, how this translates into changes in political views. These shifts in policy views help us understand which of the underlying considerations are primarily driving policy views. \\[-1.5em]
\begin{center}
\textbf{Related Literature}
\end{center} 
First, this project relates to the broad literature on measuring people's support for policies using experimental surveys. Most closely related, \textcite{stantcheva2021understanding} elicits preferences on tax policy in general. While she concentrates on the question which tax schedule of income and wealth taxes people prefer, we try to understand people's views on which instruments tax administration should use to enforce that tax schedule. 

Second, this project relates to the literature on tax enforcement, tax evasion, and especially to the literature on the role of third-party reporting in reducing tax evasion. While \textcite{kleven2011unwilling}, \textcite{pomeranz2015no}, \textcite{slemrod2017does} and \textcite{boas2024taxing} show that automatic exchange of information and third-party reporting is effective in reducing tax evasion, our research tries to uncover why many countries' parliaments and societies oppose further expanding third party reporting into their tax systems. 

Finally, our project contributes to the literature on tax filing costs and the governments' role in providing services to reduce them. \textcite{benzarti2020taxing} observes how taxpayers choose between itemizing deductions and claiming the standard deduction and finds that taxpayers forgo large tax savings to avoid compliance costs. \textcite{hauck2024optional} show that in Germany non-filing of tax declaration is common. In particular, this is the case at the lower end of the income distribution. Therefore, non-filing reduces the effectiveness of redistribution. And to show that tax compliance costs have been steadily increasing, \textcite{benzarti2024rising} use word counts of the tax codes in several countries dating back to the early 1980s as a proxy for tax compliance costs. Using a survey on US taxpayers, they also show that the majority would be willing to pay for simplifying the tax system and adopting prepopulated tax returns. Our survey expands this branch of research in several ways. First, it also looks at people's demand for tax software. Second, it elicits people's views about high levels of third-party reporting, which is in itself a prerequisite for government-provided prepopulated tax returns. Finally, by experimentally providing information, we can understand the role of (mis)information in shaping policy preferences. 
\newpage
\printbibliography
\end{document}